\documentclass[letterpaper,11pt]{Proposal}
%DECLARACIONES DE AUTORÍA
\title{Acá va el título de la propuesta}
\author{Nombre del proponente}
\date{Fecha}

%INICIA EL DOCUMENTO
\begin{document}

\begin{textblock}{10}(3,1)

\textit{Universidad de Costa Rica}\\
\textit{Facultad de Ingeniería - Escuela de Ingeniería Mecánica}\\
\textit{Anteproyecto de Maestría en Ingeniería Mecánica}

\end{textblock}

\maketitle

\begin{abstract}
Redactar acá el resumen de la propuesta, que debe concentrarse en el contenido de ese documento en particular, por lo que no debe ser de más de tres o cuatro oraciones (unas 60 palabras en total) que describan el propósito del trabajo, la metodología general y los resultados esperados.

\end{abstract}

\section{Descripción general}

Esta sección introductoria es una versión más extendida del resumen, por cuanto toca los mismos aspectos pero con un poco más de detalle, el ((suficiente)) como para que las demás secciones puedan comprenderse. Con tres o cuatro párrafos (cerca de 300 palabras) es suficiente para proveer un marco de referencia adecuado en cuanto al propósito del trabajo, la metodología general y los resultados esperados. Se pueden mencionar las empresas o instituciones que estarán directamente relacionadas con la propuesta y las posibilidades de patrocinio o financiamiento.


\section{Justificación}
En esta sección se debe explicar, en forma concisa, por qué se considera pertinente y necesario llevar a cabo el trabajo que se propone. Si bien su relación es muy cercana, no debe confundirse con la sección de antecedentes, en donde se mencionan los factores temporales y eventos que llevaron a la situación que justifica la propuesta. Se espera una extensión de 300-500 palabras.

\section{Antecedentes}
El apartado de antecedentes tiene el propósito de aclararle al lector los eventos y situaciones previas que llevaron al estado actual de necesidades que justificaron el trabajo. En medida de ello, el proponente debe conocer bien esas condiciones locales, que difícilmente se encontrarán documentadas en monografías o artículos técnicos. Se espera una extensión de 300-500 palabras.

\section{Alcance y limitaciones}

En este apartado hay que mencionar según se considere necesario las fronteras de trabajo en su extensión y profundidad técnica así como la legal. Si hay aspectos de confidencialidad, exclusión de documentos de especificaciones o planos, o exención de responsabilidad, este es el lugar en que se hace explícito. Se espera una extensión de 150 - 200 palabras.

\section{Áreas de pertinencia temática (Opcional)}

En ocasiones se estila una corta sección en donde se describen los temas generales con las que se relaciona directamente el trabajo. Típicamente se trata de mencionar de tres a cinco áreas, por lo general afines a la carrera, para orientar rápidamente al lector a la temática dominante.

\section{Marco temático de referencia}

Esta sección viene a ser una versión preliminar del marco teórico por cuanto debe mencionar el estado actual del conocimiento de los temas que son pertinentes a los objetivos del trabajo. Se espera que tenga una extensión de unas 1500 palabras. Su función es dar garantía de que se trata de un tema actual y pertinente, así como precisar la designación atinada del panel asesor.

\section{Objetivo general}

El objetivo general describe en una única oración el concepto central de la propuesta, preferiblemente sin necesitar explicaciones adicionales para un experto en el área. Una forma convencional y cómoda para redactarlas es estructurarlo en tres partes: el qué se va a hacer, cómo se va a hacer y para qué se va a hacer.

\section{Objetivos específicos}

Los objetivos específicos se pueden agrupar en dos categorías: aquellos que son parte de o contribuyen directamente al objetivo general, y los que son necesarios para el trabajo pero sólo se relacionan con el objetivo general a través del tema que se está estudiando. Precisamente por su carácter contributivo, es fácil cometer el error de redactarlos como si fueran pasos concretos para lograr el objetivo general, pero esto en realidad sería la metodología. Los objetivos específicos son hitos, no labores especificas, que se necesitan ir completando para que se pueda cumplir con el objetivo general. Se esperan de tres a cinco objetivos específicos. 

\section{Metodología}

En la propuesta de trabajo final de graduación la metodología se describe como dos o tres bloques de unos cinco a diez pasos cada uno, pero lo ideal es que está tan conciso como sea posible en consideración que es prematura establecer una metodología definitiva. Se esperan unas dos o tres viñetas de no más de 50 palabras cada una por cada objetivo especifico, pero no tiene que ser necesariamente unívoco.

\section{Patrocinio académico}

Hay que describir brevemente en esta sección el nombre, razón social, actividades generales y ubicación de la organización que está interesada o se beneficiaría directamente de los resultados y conclusiones del trabajo (300-500 palabras). También se incluye el nombre de la persona que la representará para efectos de ser la contraparte responsable de la investigación. Esta contraparte es necesaria para que obtener información al interno de la organización, coordinar autorización para las visitas pertinentes y en general ser el contacto interno. Muchas veces esta persona también estará encargada de evaluar el resultado cuando se complete el trabajo para conocer si se encuentran satisfechos con los resultados.
También se propone al panel asesor, conformado por un guía y dos lectores (de los cuales al menos uno debe ser docente activo de la Escuela de Ingeniería Mecánica). Todos los miembros del panel asesor deben tener título de Maestría o superior.

\section{Cronograma de trabajo}

Es un cuadro con la distribución general de las actividades previstas para completar los objetivos en el transcurso de los dos semestres asociados a los cursos de Taller de Investigación I (Propuesta y elaboración de marco teórico), Taller de Investigación II (ejecución y documentación metodológica) y Proyecto (Análisis de resultados, informe final y defensa formal). Se recomienda mencionar los elementos descritos en la metodología.

\begin{table}[!hbt]
\begin{tabularx}{\columnwidth}{p{93mm} *{12}{W}}		
	\noalign{\hrule height 1.5pt}
    \centering \multirow{2}{*}{Actividad}&\multicolumn{12}{c}{Mes}\\
    \cline{2-13}
    &1&2&3&4&5&6&7&8&9&10&11&12\\
    \hline
    1. Objetivo 1.&\cellcolor{black}&&&&&&&&&&&\\
    2. Objetivo 2.&&\multicolumn{2}{c}{\cellcolor{black}}&&&&&&&&&\\
    3. Objetivo 3.&&&&\multicolumn{3}{c}{\cellcolor{black}}&&&&&&\\
    4. Objetivo 4.&&&&&&&\multicolumn{4}{c}{\cellcolor{black}}&&\\
 	\noalign{\hrule height 1.5pt}         \end{tabularx}
            \end{table}
\section{Estructura de contenidos}


Este borrador de índice es la versión prevista de los nombres de los capítulos y sus secciones, aún a sabiendas de que está muy distante de completarse. Se recomienda que se presente dentro de un cuadro de texto para que no se confunda con las secciones de la propuesta ni se altere su enumeración. Dado que recibirá muchos cambios a lo largo de la ejecución de la propuesta, no debe considerarse vinculante; es una referencia para describir el entregable del proyecto de graduación.

\section{Referencias bibliográficas}

Toda información citada deberá llevar la referencia bibliográfica correspondiente, que se incorporan en esta sección junto con aquellos artículos o monografías que se consultaron antes y se usaron al elaborar la propuesta.

%\bibliographystyle{plain}
%\bibliography{misbilos}

\end{document}